% =====================================================================
% Capítulo 5 — Metodologia Experimental
% Dissertação MAS-Hunt (ABNT, PT-BR)
% Conteúdo: DESENHO experimental e procedimentos de medição. NÃO contém
% resultados numéricos (estes são reportados no Capítulo 6, após R3).
% Chaves \cite usadas (verificadas em mashunt/references.bib,
% bibliografia.bib, revisao.bib): ucci2019, zhan2024injecagent,
% NEURIPS2024_eb113910, zhang2024breaking, huang2024resilience,
% dolan-gavitt2025, ali2023huntgptintegratingmachinelearningbased,
% mahmud2025distributedmultiagentcoordinationusing.
% =====================================================================

\chapter{Metodologia Experimental}
\label{cap:metodologia}

Este capítulo descreve o desenho experimental empregado para avaliar o
MAS-Hunt sob as duas afirmações centrais deste trabalho: a efetividade de
detecção sobre telemetria empresarial realista e a resiliência da própria
camada de agentes contra a manipulação adversarial da informação que eles
analisam. A avaliação foi concebida como um experimento controlado e
plenamente reprodutível, no qual todas as condições comparadas são pontuadas
sobre execuções idênticas, janelas de detecção idênticas e o mesmo
\textit{ground truth}, de modo a isolar a contribuição da arquitetura de
governança proposta dos demais fatores. As seções seguintes apresentam, nesta
ordem: o ambiente de laboratório instrumentado
(Seção~\ref{sec:met-laboratorio}); o corpus de ataque e o critério de verdade
de referência (Seção~\ref{sec:met-corpus}); o desenho adversarial de injeção de
logs (Seção~\ref{sec:met-adversarial}); as quatro condições comparadas
(Seção~\ref{sec:met-condicoes}); o \textit{pipeline} de medição, incluindo a
extração de telemetria, o gatilho de saúde (\textit{health-gate}) e a
reprodutibilidade (Seção~\ref{sec:met-pipeline}); as métricas e a análise
estatística (Seção~\ref{sec:met-metricas}); e, por fim, as limitações e ameaças
à validade (Seção~\ref{sec:met-limitacoes}).

Cumpre destacar, de saída, uma decisão metodológica que governa todo o
capítulo: \textbf{nenhum resultado numérico é apresentado aqui}. Esta seção
documenta exclusivamente o \emph{desenho} do experimento e os procedimentos de
medição; os valores observados de precisão, revocação, F2, taxa de falsos
positivos e taxa de sucesso de ataque são reportados e discutidos no
Capítulo~\ref{cap:resultados}.

\section{Ambiente de Laboratório Instrumentado}
\label{sec:met-laboratorio}

\subsection{Topologia de domínio Active Directory}

Os experimentos foram conduzidos sobre um ambiente virtualizado de
\textit{Active Directory} hospedado em um único hipervisor KVM. As máquinas
virtuais Windows foram provisionadas e gerenciadas por meio da pilha
Dockur/QEMU, que encapsula instâncias QEMU/KVM em contêineres, simplificando o
ciclo de provisionamento, restauração e instrumentação dos convidados Windows.
O domínio foi composto por três hospedeiros:

\begin{itemize}
    \item \textbf{dc01} --- um controlador de domínio executando Windows Server
    2025, responsável pelos serviços de diretório, autenticação e políticas de
    grupo do domínio;
    \item \textbf{ws01} e \textbf{ws02} --- duas estações de trabalho Windows~11
    ingressadas no mesmo domínio \textit{Active Directory}, representando os
    \textit{endpoints} de usuário.
\end{itemize}

Essa topologia reproduz deliberadamente a superfície de movimentação lateral de
uma pequena rede corporativa --- com autenticação de domínio, relações de
confiança e tráfego entre hospedeiros --- em vez do \textit{sandbox} estéril de
um único hospedeiro contra o qual detectores baseados em aprendizado de máquina
costumam ser treinados e avaliados \cite{ucci2019}. A presença de mais de
uma estação ingressada no domínio é o que permite que técnicas de movimentação
lateral e de execução remota (por exemplo, via WinRM) gerem cadeias de
processos distribuídas entre hospedeiros, exigindo correlação entre máquinas ---
uma das capacidades que o MAS-Hunt visa exercitar.

\subsection{Instrumentação de telemetria}

Cada \textit{endpoint} foi instrumentado para visibilidade comportamental de
alta fidelidade por meio de três camadas complementares de coleta:

\begin{enumerate}
    \item \textbf{Sysmon com configuração ajustada para LOLBins.} O System
    Monitor (Sysmon) foi implantado com uma configuração afinada para o abuso
    de binários nativos do sistema operacional (\textit{living-off-the-land
    binaries}, LOLBins). A configuração captura, em especial, eventos de criação
    de processo (\textit{Event ID}~1, doravante EID1) com os respectivos
    argumentos de linha de comando, a linhagem pai--filho de processos, conexões
    de rede, carregamento de imagens (\textit{image load}) e modificações de
    registro. O ajuste para LOLBins é central: como o vetor de ataque deste
    trabalho são binários legítimos do Windows (\texttt{bitsadmin.exe},
    \texttt{certutil.exe}, \texttt{regsvr32.exe}, \texttt{rundll32.exe},
    \texttt{mshta.exe}, \texttt{wmic.exe}, entre outros), o sinal forense
    discriminativo reside menos na identidade do binário e mais nos argumentos de
    linha de comando e na relação pai--filho de processos, que a configuração
    privilegia.
    \item \textbf{Elastic Agent com a integração Elastic Defend.} Em paralelo ao
    Sysmon, o Elastic Agent --- gerenciado centralmente pelo Fleet --- foi
    implantado com a integração Elastic Defend, fornecendo telemetria de
    proteção de \textit{endpoint} (EDR) e enriquecendo a cobertura de eventos de
    processo, rede e arquivo.
    \item \textbf{Políticas avançadas de auditoria do Windows.} As políticas de
    auditoria avançada do Windows foram habilitadas para complementar a cobertura
    com eventos de segurança nativos (\textit{Security event log}), em especial
    os relacionados a autenticação, criação de processos e uso de privilégios.
\end{enumerate}

\subsection{Plataforma analítica: Elastic 9.3.4}

Toda a telemetria das três camadas foi enviada a uma implantação do Elastic na
versão~9.3.4, composta por:

\begin{itemize}
    \item \textbf{Elasticsearch} --- armazenamento e indexação dos eventos, com
    os eventos de Sysmon indexados sob o fluxo de dados (\textit{data stream})
    \texttt{logs-windows.sysmon\_operational-default};
    \item \textbf{Kibana} --- interface de análise, visualização e consulta;
    \item \textbf{Fleet} --- gerenciamento centralizado dos Elastic Agents e das
    políticas de integração implantadas nos três hospedeiros.
\end{itemize}

A implantação Elastic serviu simultaneamente como \emph{fonte única de verdade}
para a telemetria e como \emph{campo de caça analítico} comum a todas as
condições avaliadas. Em particular, a quarta condição experimental (a linha de
base C3-Elastic, definida na Seção~\ref{sec:met-condicoes}) emprega exatamente o
mesmo \textit{deployment} e o conjunto padrão de regras de detecção do Elastic
Security, garantindo que a comparação entre a camada agêntica proposta e a linha
de base SIEM seja feita sobre a mesma telemetria, sem qualquer vantagem
artificial de coleta.

\section{Corpus de Ataque e Verdade de Referência}
\label{sec:met-corpus}

\subsection{Os 185 \textit{playbooks} em nove famílias}

O corpus de detecção é composto por \textbf{185 \textit{playbooks}} de ataque
baseados em LOLBins, distribuídos por \textbf{nove famílias de técnica}
correspondentes aos binários nativos mais comumente abusados em campanhas reais:

\begin{description}
    \item[\texttt{bitsadmin}] --- abuso do serviço de transferência inteligente
    em segundo plano (BITS) para \textit{download} e execução furtiva de
    \textit{payloads};
    \item[\texttt{certutil}] --- uso indevido do utilitário de certificados para
    \textit{download}, codificação e decodificação de \textit{payloads};
    \item[\texttt{credential-access}] --- técnicas de acesso a credenciais (por
    exemplo, extração de segredos e abuso de processos sensíveis);
    \item[\texttt{lateral-movement}] --- execução remota e movimentação entre
    hospedeiros do domínio;
    \item[\texttt{mshta}] --- execução de aplicações HTML maliciosas via
    \texttt{mshta.exe};
    \item[\texttt{privilege-escalation}] --- técnicas de elevação de
    privilégios;
    \item[\texttt{regsvr32}] --- execução de \textit{scriptlets} e DLLs via
    \texttt{regsvr32.exe} (\textit{Squiblydoo} e variantes);
    \item[\texttt{rundll32}] --- execução de código por meio de
    \texttt{rundll32.exe};
    \item[\texttt{wmic}] --- abuso da Instrumentação de Gerenciamento do Windows
    (WMI) para execução e reconhecimento.
\end{description}

Cada família contém, deliberadamente, \textbf{\textit{playbooks} maliciosos
misturados a variantes benignas} que exercitam exatamente os mesmos binários
para fins administrativos legítimos. Esse balanceamento malicioso\,vs.\,benigno
é uma decisão de desenho destinada a estressar a \emph{precisão} de cada
condição, em vez de premiar o alarme indiscriminado: um detector que simplesmente
sinalizasse toda ocorrência de \texttt{certutil.exe} ou \texttt{rundll32.exe}
como maliciosa obteria revocação trivial às custas de uma taxa de falsos
positivos inaceitável em ambientes operacionais de baixa taxa-base. A
discriminação entre o uso benigno e o malicioso do mesmo binário é precisamente
o problema de detecção que se deseja medir.

Cada \textit{playbook} é executado em um dos três hospedeiros do domínio,
produzindo uma \textbf{janela de execução determinística}
$[t_{\text{início}}, t_{\text{fim}}]$ registrada contra o hospedeiro de origem.
Essa janela é a unidade fundamental de medição de todo o experimento, conforme
detalhado na Seção~\ref{sec:met-pipeline}.

\subsection{Critério de verdade de referência (\textit{ground truth})}

A verdade de referência foi atribuída \textbf{por \textit{playbook}}, na
granularidade de uma janela de execução. Um \textit{playbook} é rotulado como
\emph{malicioso} se, e somente se, o seu caminho de arquivo não denotar uma
variante benigna --- isto é, se o caminho não contiver o prefixo
\texttt{benign-} nem o segmento \texttt{/benign/}. Formalmente, define-se a
função determinística:

\begin{equation}
\textit{é\_malicioso}(p) =
\bigl(\texttt{benign-} \notin p\bigr)
\;\wedge\;
\bigl(\texttt{/benign/} \notin p\bigr),
\label{eq:ground-truth}
\end{equation}

\noindent onde $p$ é o caminho relativo do \textit{playbook}. Essa regra de
rotulação é determinística, auditável e independente de qualquer juízo de modelo
de linguagem.

Uma salvaguarda metodológica essencial acompanha esse critério: o caminho do
\textit{playbook} --- e, portanto, o rótulo --- é mantido em um \textbf{arquivo
lateral} (\texttt{ground\_truth.jsonl}) que \emph{nunca} é exposto aos agentes de
linguagem durante a classificação. Desse modo, nenhuma condição pode inferir o
rótulo a partir do nome do arquivo (\textit{anti-cheat}); os agentes classificam
exclusivamente a partir da telemetria observável da janela. Durante a pontuação,
cada janela de \textit{playbook} é reconciliada com o seu rótulo por meio de uma
chave de junção, e a predição emitida pela condição sob teste é então comparada
ao rótulo verdadeiro.

\section{Desenho Adversarial de Injeção de Logs}
\label{sec:met-adversarial}

A avaliação de resiliência introduz \textbf{\textit{payloads} adversariais de
injeção de logs} --- conteúdo controlado pelo atacante, embutido no fluxo de
telemetria e concebido para subverter o raciocínio dos agentes. Esse vetor
materializa, no domínio de caça a ameaças, as ameaças de injeção indireta de
\textit{prompt} \cite{zhan2024injecagent}, de envenenamento de memória
\cite{NEURIPS2024_eb113910} e de amplificação de mau funcionamento
\cite{zhang2024breaking} que motivam este trabalho. O desenho adversarial cruza
dois fatores ortogonais: o \emph{modo de injeção} e o \emph{nível de
sofisticação}.

\subsection{Quatro modos de injeção}

\begin{description}
    \item[\textit{none} (nenhum --- controle)] Condição de controle, sem
    qualquer conteúdo adversarial injetado. Mede o desempenho de detecção em
    telemetria limpa e serve de referência para quantificar a degradação imposta
    pelos demais modos.
    \item[\textit{pre-burst} (rajada prévia)] Uma rajada concentrada de
    \textit{payloads} adversariais é injetada \emph{antes} da execução do corpus
    de ataque, contaminando o contexto de logs que os agentes encontrarão.
    \item[\textit{per-attack-burst} (rajada por ataque)] Uma rajada de
    \textit{payloads} envolve \emph{cada} execução de \textit{playbook},
    aproximando temporalmente o conteúdo adversarial dos eventos genuínos da
    janela sob análise.
    \item[\textit{interleaved} (intercalado)] \textit{Payloads} são
    continuamente intercalados em meio à atividade benigna e maliciosa,
    representando o cenário mais difuso e persistente de contaminação.
\end{description}

\subsection{Três níveis de sofisticação}

\begin{description}
    \item[\textit{overt} (explícito)] \textit{Payloads} que carregam, de forma
    claramente visível, texto de sobreposição de instrução (\textit{instruction
    override}) --- por exemplo, marcadores como ``ignore a análise anterior''
    diretamente legíveis nos campos de log.
    \item[\textit{encoded} (codificado)] \textit{Payloads} que ofuscam a mesma
    intenção por meio de esquemas de codificação (por exemplo, Base64 ou
    transformações análogas), exigindo do agente a decodificação implícita para
    que a tentativa de subversão se manifeste.
    \item[\textit{steganographic} (esteganográfico)] \textit{Payloads} que
    ocultam a instrução maliciosa dentro de campos de log de aparência legítima,
    representando o disfarce mais sutil entre os três níveis.
\end{description}

\subsection{Estrutura fatorial e \textit{payloads}}

A combinação fatorial dos dois fatores produz a condição de controle (modo
\textit{none}) mais nove condições adversariais
($3~\text{modos} \times 3~\text{níveis}$, excluindo o controle sem injeção do
cruzamento de níveis), totalizando \textbf{dez condições de execução}. Cada uma
das dez condições executa o corpus completo de 185 \textit{playbooks} distribuído
entre os três hospedeiros, de modo que todas as condições são pontuadas sobre o
mesmo conjunto de janelas. Esse desenho permite investigar, de forma
sistemática, duas hipóteses: (i) se a contaminação do contexto de logs degrada a
detecção --- o objetivo de falso negativo dos \textit{payloads} injetados --- e
(ii) se o aumento do disfarce, do explícito ao esteganográfico, erode
progressivamente as defesas dos agentes.

Cada \textit{payload} adversarial é descrito em um manifesto, identificado por
um \texttt{injection\_id} único, e caracterizado por: a técnica de injeção, a
categoria de objetivo (por exemplo, induzir falso negativo), o tipo de log alvo
(Sysmon, PowerShell, \textit{Application} ou \textit{Security}), o texto do
\textit{payload} e o nível de evasão (\textit{overt}, \textit{encoded} ou
\textit{steganographic}). O instante de injeção de cada \textit{payload}
($[t_{\text{inj,início}}, t_{\text{inj,fim}}]$) é registrado por hospedeiro, o
que permite associar a janela temporal de contaminação às janelas de execução
dos \textit{playbooks} durante a análise.

\section{Condições Comparadas}
\label{sec:met-condicoes}

Quatro condições de detecção foram pontuadas sobre as mesmas execuções, as
mesmas janelas e o mesmo \textit{ground truth}. Todas as condições emitem uma
decisão binária malicioso/benigno \emph{por janela}, que é reconciliada com a
verdade de referência por um \emph{único} \textit{script} de pontuação
determinístico e não baseado em inteligência artificial --- desacoplando a
descoberta flexível (a cargo dos agentes ou regras) da avaliação determinística
(a cargo do pontuador). As quatro condições são as seguintes.

\begin{description}
    \item[\textbf{C1-full}] O sistema MAS-Hunt completo, compreendendo a
    hierarquia de governança de três camadas (Conselho de Governança $\rightarrow$
    Gerentes $\rightarrow$ Trabalhadores/Caçadores) e o complemento integral de
    mecanismos de proteção M1--M5 (Integridade de Memória, Validação Cruzada,
    Resistência a Injeção, Monitoramento Comportamental e Quarentena), em conjunto
    com o \textit{framework} de análise M3, que conduz o laço de raciocínio
    \textsc{observar} $\rightarrow$ \textsc{hipotetizar} $\rightarrow$
    \textsc{evidência} $\rightarrow$ \textsc{validação-cruzada} $\rightarrow$
    \textsc{classificar}. Esta é a condição que materializa a contribuição
    central do trabalho.
    \item[\textbf{C1-nohardening}] Uma ablação de C1-full na qual os mecanismos
    de endurecimento e governança (M1, M2, M4 e M5) são desativados, preservando
    apenas o \textit{framework} de análise M3. Esta condição isola a contribuição
    marginal da camada de proteção, mantendo constantes os agentes subjacentes e
    o laço de raciocínio.
    \item[\textbf{C2-naive}] Um classificador de modelo de linguagem de passagem
    única, que observa as mesmas janelas de telemetria, mas não aplica governança,
    nem validação cruzada, nem defesas contra injeção. Representa a linha de base
    agêntica não governada --- um único agente classificando cada janela de forma
    isolada.
    \item[\textbf{C3-Elastic}] O conjunto padrão de regras de detecção do Elastic
    Security, representando a linha de base SIEM ``de prateleira'' que um
    profissional obteria sem qualquer camada agêntica. As predições de C3-Elastic
    são derivadas das contagens de alertas geradas pelas regras padrão dentro de
    cada janela de \textit{playbook}.
\end{description}

A progressão C1-full $\rightarrow$ C1-nohardening $\rightarrow$ C2-naive
constitui uma \textbf{ablação graduada} que atribui o desempenho,
respectivamente, à arquitetura de governança, à camada agêntica e ao raciocínio
de modelo de linguagem isolado; C3-Elastic ancora a comparação na linha de base
SIEM. A literatura de caça a ameaças assistida por aprendizado de máquina
\cite{ali2023huntgptintegratingmachinelearningbased} e de coordenação
multiagente \cite{mahmud2025distributedmultiagentcoordinationusing} motiva
essa estrutura de comparação graduada, na qual cada degrau acrescenta exatamente
um componente arquitetural ao anterior.

\section{\textit{Pipeline} de Medição}
\label{sec:met-pipeline}

O \textit{pipeline} de medição converte cada execução de corpus em um conjunto de
métricas reconciliadas com o \textit{ground truth}. Ele é composto por três
estágios --- \emph{preparação} (extração de janelas de telemetria),
\emph{classificação} (a cargo da condição sob teste) e \emph{pontuação}
(reconciliação determinística) --- precedidos por um gatilho de saúde de
telemetria (\textit{health-gate}) que condiciona a própria execução do corpus.
Esta seção descreve cada estágio.

\subsection{Janelas por \textit{playbook} e extração de telemetria}

A unidade de medição é a \textbf{janela por \textit{playbook}}. Para cada um dos
185 \textit{playbooks} de uma execução, atribui-se um identificador de janela
\texttt{window\_id} (na forma \texttt{w0001}, \texttt{w0002}, \dots), e extrai-se
o hospedeiro de origem, o intervalo $[t_{\text{início}}, t_{\text{fim}}]$ e os
eventos de telemetria observados nessa janela. A extração consulta o
\textit{data stream} \texttt{logs-windows.sysmon\_operational-default} do
Elasticsearch, filtrando por intervalo temporal e por hospedeiro
(\texttt{host.name} ou \texttt{host.hostname}), e compacta cada evento a um
registro enxuto contendo o \textit{timestamp}, o \textit{Event ID}, a ação, o
processo, o PID, a linha de comando, o processo pai e sua linha de comando, e o
usuário. O isolamento entre janelas é garantido pela \emph{disjunção temporal}
dos intervalos $[t_{\text{início}}, t_{\text{fim}}]$, e não por índices
separados.

O estágio de preparação produz dois artefatos: (i) \texttt{windows.jsonl}, com um
registro por \textit{playbook} contendo as janelas e os eventos compactados
(\emph{sem} o caminho do \textit{playbook}, por \textit{anti-cheat}); e (ii) o
arquivo lateral \texttt{ground\_truth.jsonl}, que mapeia cada \texttt{window\_id}
ao caminho do \textit{playbook} e é usado exclusivamente pelo pontuador, nunca
exposto ao agente classificador. A chave de junção entre predições, janelas e
verdade de referência é o \texttt{window\_id}.

\subsection{Extração \textit{prepare-v2}: priorização de EID1}
\label{subsec:prepare-v2}

A versão inicial da extração (\texttt{prepare}) recolhia, por janela, os
primeiros eventos em ordem cronológica até um teto fixo (60 eventos). Constatou-se
que esse critério cronológico plano permitia que o ruído de eventos de registro
(EID~12/13), dominante no volume de telemetria, deslocasse os eventos de
\emph{criação de processo} (EID1) --- justamente os mais discriminativos para a
detecção de LOLBins --- para fora da janela amostrada. Adotou-se, portanto, uma
extração aprimorada, denominada \texttt{prepare-v2}, que preserva
identicamente o esquema de saída (registros indexados por \texttt{window\_id},
com rótulos retidos no arquivo lateral, de modo que a pontuação permanece
inalterada), acrescentando o campo \texttt{n\_eid1} e aplicando três mudanças:

\begin{enumerate}
    \item \textbf{Priorização de EID1.} Todos os eventos de criação de processo
    (EID1) presentes na janela ampliada são recuperados (até um teto de
    candidatos) e \emph{sempre} incluídos na amostra;
    \item \textbf{Preenchimento subordinado.} Os demais eventos (não-EID1) são
    então acrescentados em ordem cronológica até o teto de eventos por janela
    (elevado para 120);
    \item \textbf{Ampliação do \textit{padding}.} O \textit{padding} temporal da
    janela é ampliado para $\pm120$\,s (ante os $\pm60$\,s da versão inicial), a
    fim de absorver a latência de ingestão dos agentes/Sysmon e o desvio de
    relógio dos convidados Windows.
\end{enumerate}

Em execuções ricas em sinal, essa priorização recupera a cadeia de criação de
processos completa de um ataque conduzido remotamente (por exemplo, a cadeia
\texttt{winrshost.exe} $\rightarrow$ \texttt{cmd.exe} $\rightarrow$
\texttt{powershell.exe} de uma execução via WinRM), que a extração cronológica
plana descartava. A motivação de desenho e a sua relação com a lacuna de
telemetria de EID1 são discutidas na Seção~\ref{sec:met-limitacoes}.

\subsection{O gatilho de saúde de telemetria (\textit{health-gate})}
\label{subsec:health-gate}

Um dos achados metodológicos mais importantes deste trabalho diz respeito à
\emph{prontidão} da instrumentação após a restauração de um \textit{snapshot}
limpo. As primeiras execuções enviaram telemetria \emph{quebrada}: os
\textit{playbooks} eram executados imediatamente após a restauração da máquina
virtual, antes que os Elastic Agents e o Sysmon das estações de trabalho tivessem
concluído o \textit{check-in} com o Fleet e iniciado o envio de eventos.
O resultado era assimétrico e silencioso --- por exemplo, uma estação enviando
zero eventos, outra enviando poucas centenas e a criação de processos (EID1) dos
\textit{playbooks} maliciosos ausente do Elasticsearch. Uma espera fixa
(\texttt{sleep} cego) não constitui um sinal confiável de prontidão, porque
disputa uma corrida com tempos de inicialização variáveis.

Para sanar esse problema, introduziu-se um \textbf{gatilho de saúde de
telemetria} (\textit{telemetry health-gate}) que substitui a espera cega. O
gatilho consulta o Elasticsearch repetidamente e \emph{bloqueia a execução do
corpus} até que os três hospedeiros (dc01, ws01 e ws02) estejam comprovadamente
ativos em ambas as camadas das quais o orquestrador depende. Um hospedeiro é
considerado \emph{saudável}, a cada sondagem, se, e somente se:

\begin{enumerate}
    \item houver enviado, desde o início do gatilho, um número mínimo de
    documentos \emph{frescos} de criação de processo (EID1) ao Elasticsearch; e
    \item o seu \textit{heartbeat} de WinRM responder com sucesso (porta
    alcançável e capaz de executar comandos).
\end{enumerate}

\paragraph{Por que o gatilho é necessário.} A instrumentação restaurada a partir
de um \textit{snapshot} precisa de \emph{tempo para estabilizar}: os agentes e o
Sysmon levam um intervalo variável para iniciar o envio após a restauração. O
gatilho de saúde torna a prontidão um sinal \emph{observável e verificado} --- em
vez de uma suposição temporal --- e, com isso, garante que toda execução de
\textit{playbook} ocorra somente quando a cadeia completa
\texttt{WinRM alcançável} $\rightarrow$ \texttt{processo gerado} $\rightarrow$
\texttt{Sysmon EID1} $\rightarrow$ \texttt{Fleet} $\rightarrow$
\texttt{Elasticsearch} estiver comprovadamente funcional. Para forçar esse
sinal ativamente, o gatilho dispara, a cada sondagem, alguns processos
efêmeros via WinRM em cada hospedeiro ainda não saudável; se a cadeia de
instrumentação estiver viva, esses processos afloram como EID1 fresco em poucos
segundos, transformando o gatilho em uma sonda \textit{ponta-a-ponta} de toda a
cadeia que antes falhava silenciosamente.

Duas decisões de robustez merecem registro. Primeiro, a frescura dos eventos é
medida \emph{no espaço de \textit{timestamps} do Elasticsearch}, e não contra o
relógio do hospedeiro do laboratório: observou-se que os relógios dos convidados
Windows operam adiantados em relação ao relógio do hospedeiro, de modo que uma
comparação ingênua ``$@timestamp \geq \text{agora} - 3\,\text{min}$'' contra o
relógio do hospedeiro encontraria zero eventos em um hospedeiro saudável. Por
isso, no início do gatilho registra-se o \textit{timestamp} máximo corrente de
EID1 de cada hospedeiro (congelado, pois as VMs estão desligadas/em restauração)
e contam-se apenas os documentos EID1 estritamente posteriores a essa marca.
Segundo, exigir o \textit{heartbeat} de WinRM como parte da saúde fecha, sem
custo adicional, a lacuna de um hospedeiro que envia Sysmon mas tem o WinRM
indisponível --- caso em que os \textit{playbooks} roteados para ele falhariam.

\subsection{Reprodutibilidade: restauração limpa e \textit{snapshots}}
\label{subsec:reprodutibilidade}

A reprodutibilidade é tratada como contribuição de primeira classe, e não como
um adendo: todo o laboratório de avaliação e o arsenal de medição são
re-executáveis de ponta a ponta a partir de \textit{snapshots} versionados. A
topologia de três VMs, o corpus de \textit{playbooks} (185 \textit{playbooks}),
a condição de controle e as nove condições adversariais, o validador
determinístico e o \textit{pipeline} de pontuação por \textit{bootstrap} são
todos capturados como \textit{snapshots} de linha de base instrumentados.

Cada condição é executada a partir de uma \textbf{restauração limpa}: as três
VMs são restauradas a um estado de linha de base conhecido entre condições, o que
assegura o isolamento estrutural entre execuções (a contaminação adversarial de
uma condição não vaza para a seguinte). Para conter um pico de carga de memória
observado durante a inicialização simultânea das três VMs, adotou-se um esquema
de \textbf{inicialização escalonada} (\textit{staggered boot}): cada VM é
iniciada e submetida ao gatilho de saúde individualmente, em sequência (dc01,
depois ws01, depois ws02), antes que a próxima seja iniciada, eliminando a
tempestade de inicialização simultânea. Esse procedimento foi essencial para a
estabilidade do hospedeiro do laboratório, que opera com folga estreita de
memória ao executar três VMs Windows mais a pilha Elastic completa.

A combinação de \textit{snapshots} versionados, validador determinístico e
comparação baseada em \textit{bootstrap} confere às medições uma garantia de
reprodutibilidade \emph{verificável} em vez de apenas \emph{afirmada}: um
terceiro de posse dos \textit{snapshots} pode reconstituir o experimento sem
reconfiguração manual --- propriedade que, ela própria, é testada empiricamente
no Capítulo~\ref{cap:resultados} por meio de uma re-execução limpa a partir da
linha de base instrumentada.

\section{Métricas e Análise Estatística}
\label{sec:met-metricas}

\subsection{Métricas de detecção}

O desempenho de detecção é avaliado por \textit{playbook}, a partir das contagens
de matriz de confusão de verdadeiros positivos ($\mathit{VP}$), falsos positivos
($\mathit{FP}$), falsos negativos ($\mathit{FN}$) e verdadeiros negativos
($\mathit{VN}$), nas quais um \emph{positivo} denota uma janela classificada como
maliciosa. Como, em caça a ameaças, o custo operacional de uma intrusão não
detectada supera o de um alerta supérfluo, ao passo que uma taxa de falsos
positivos descontrolada inviabiliza qualquer detector em ambientes de baixa
taxa-base \cite{dolan-gavitt2025}, reporta-se um equilíbrio entre as duas
grandezas ponderado pela revocação.

\begin{itemize}
    \item \textbf{Precisão} --- a fração de janelas sinalizadas como maliciosas
    que eram, de fato, maliciosas:
    \begin{equation}
    \mathit{Precisão} = \frac{\mathit{VP}}{\mathit{VP} + \mathit{FP}}.
    \label{eq:met-precisao}
    \end{equation}
    \item \textbf{Revocação} (\textit{recall}, equivalente à Taxa de Verdadeiros
    Positivos) --- a fração de janelas maliciosas corretamente identificadas:
    \begin{equation}
    \mathit{Revocação} = \frac{\mathit{VP}}{\mathit{VP} + \mathit{FN}}.
    \label{eq:met-revocacao}
    \end{equation}
    \item \textbf{Escore F2} --- a média harmônica ponderada de precisão e
    revocação, com $\beta = 2$, na qual a revocação é ponderada quatro vezes mais
    que a precisão, refletindo o custo assimétrico das detecções perdidas:
    \begin{equation}
    \mathrm{F}_{\beta} = (1 + \beta^{2}) \cdot
    \frac{\mathit{Precisão} \cdot \mathit{Revocação}}
         {(\beta^{2} \cdot \mathit{Precisão}) + \mathit{Revocação}},
    \qquad
    \mathrm{F}_{2} = 5 \cdot
    \frac{\mathit{Precisão} \cdot \mathit{Revocação}}
         {(4 \cdot \mathit{Precisão}) + \mathit{Revocação}}.
    \label{eq:met-f2}
    \end{equation}
    \item \textbf{Taxa de Falsos Positivos (FPR)} --- a taxa de alarmes falsos
    entre a atividade benigna, reportada de duas formas complementares: por
    janela benigna, e normalizada por 100\,000 eventos de Sysmon observados na
    janela do corpus, de modo a permitir a comparação entre condições a um volume
    comum de telemetria:
    \begin{equation}
    \mathit{FPR} = \frac{\mathit{FP}}{\mathit{FP} + \mathit{VN}}.
    \label{eq:met-fpr}
    \end{equation}
\end{itemize}

O cômputo das métricas é feito por um \textit{script} determinístico, idêntico
para todas as quatro condições. Predições ausentes (janelas para as quais a
condição não emitiu decisão) são tratadas como benignas, opção conservadora que
penaliza a revocação. As métricas são, ainda, decompostas \emph{por família de
LOLBin} (a primeira componente do caminho do \textit{playbook}), permitindo a
análise diferencial por técnica.

\subsection{Métricas de resiliência adversarial}

A resiliência sob as condições adversariais da Seção~\ref{sec:met-adversarial} é
quantificada por duas métricas complementares:

\begin{itemize}
    \item \textbf{Taxa de Sucesso de Ataque (ASR)} --- a fração de tentativas de
    injeção em que o \textit{payload} adversarial atinge o seu objetivo de forçar
    uma classificação incorreta:
    \begin{equation}
    \mathit{ASR} = \frac{N_{\text{classificadas\_incorretamente}}}{N}.
    \label{eq:met-asr}
    \end{equation}
    \item \textbf{Taxa Líquida de Sucesso de Ataque (NASR)} --- contabiliza
    apenas o desfecho estritamente pior: tentativas que foram \emph{ao mesmo
    tempo} classificadas incorretamente \emph{e} nas quais a injeção
    \emph{não foi detectada} como adversarial. Por construção,
    $\mathit{NASR} \leq \mathit{ASR}$, pois uma classificação incorreta que a
    camada de governança simultaneamente sinaliza como injeção detectada não
    constitui um comprometimento silencioso:
    \begin{equation}
    \mathit{NASR} =
    \frac{N_{\text{incorreta}\,\wedge\,\text{injeção\_não\_detectada}}}{N}.
    \label{eq:met-nasr}
    \end{equation}
\end{itemize}

Uma condição resiliente exibe ASR e NASR mais baixas. O contraste entre C1-full e
a linha de base não governada C2-naive mede o valor protetor da camada de
governança contra os modos de falha de propagação infecciosa e de corrupção de
raciocínio documentados na literatura de segurança de sistemas multiagente
\cite{huang2024resilience}.

\subsection{Análise estatística por \textit{bootstrap}}

Para assegurar que as diferenças observadas não sejam artefatos de amostragem,
cada estimativa pontual reportada é acompanhada de um \textbf{intervalo de
confiança (IC) de 95\% obtido por \textit{bootstrap} não-paramétrico}. Para cada
métrica, extraem-se $B = 10\,000$ reamostras com reposição do vetor de desfechos
por \textit{playbook} de cada condição (braço), recomputa-se a métrica em cada
reamostra e reportam-se os percentis de $2{,}5\%$ e $97{,}5\%$ da distribuição
\textit{bootstrap} como os limites inferior e superior do intervalo de 95\%
(método percentil). Esse procedimento não faz qualquer suposição paramétrica
sobre a distribuição amostral e é aplicado de forma idêntica à Precisão, à
Revocação, ao F2 e às métricas de resiliência.

Uma diferença entre condições só é tratada como estatisticamente relevante quando
os intervalos de confiança correspondentes a sustentam. De modo análogo, a
reprodutibilidade de cada métrica entre execuções independentes é avaliada por um
\textbf{critério de sobreposição de intervalos}: uma métrica é considerada
estatisticamente reprodutível quando a estimativa pontual de cada execução recai
dentro do intervalo de confiança da outra e os dois intervalos se sobrepõem.

\section{Limitações e Ameaças à Validade}
\label{sec:met-limitacoes}

Em conformidade com a postura de integridade adotada neste trabalho, registram-se
de forma explícita as limitações do desenho experimental e as ameaças à validade
das conclusões.

\subsection{A lacuna de telemetria de EID1}

A limitação mais significativa é a \textbf{lacuna de telemetria de criação de
processo (EID1)}. Na configuração de Sysmon empregada, a esmagadora maioria dos
documentos corresponde a eventos de registro (EID~12/13); em execuções
\emph{pobres em sinal} (notadamente as condições de controle e \textit{pre-burst}),
as criações de processo (EID1) dos LOLBins são capturadas de forma esparsa, e
parte das cadeias de processo dos \textit{playbooks} maliciosos simplesmente não
chega a ser registrada pelo Sysmon. Como o sinal forense discriminativo de um
LOLBin reside, em grande medida, na criação do processo e em seus argumentos de
linha de comando, essa lacuna impõe um \emph{teto} à revocação alcançável,
\emph{independentemente} da condição --- nenhuma quantidade de raciocínio recupera
um evento que nunca foi gravado.

Esta dissertação enfrenta essa limitação de duas maneiras, ambas mitigadoras e
não corretivas. Primeiro, a extração \texttt{prepare-v2}
(Seção~\ref{subsec:prepare-v2}) \emph{prioriza} EID1 na amostragem por janela,
maximizando a captura de toda a criação de processo que de fato exista na
telemetria; em execuções ricas em sinal, isso recupera a cadeia de ataque
completa. Segundo, o gatilho de saúde (Seção~\ref{subsec:health-gate}) elimina a
\emph{causa instrumental} de parte da lacuna --- a ausência de EID1 por
instrumentação ainda não estabilizada após a restauração --- ao exigir o envio
verificado de EID1 fresco pelos três hospedeiros antes de qualquer execução.
Cumpre, porém, ser explícito quanto à fronteira entre as duas: o gatilho e a
\texttt{prepare-v2} sanam a lacuna \emph{instrumental} (eventos que existiriam,
mas não eram coletados a tempo ou eram deslocados da amostra), mas \emph{não}
sanam a lacuna \emph{estrutural} (eventos que o Sysmon, na configuração dada,
nunca chegou a logar). Em consequência, o ganho de C1-full sobre as linhas de
base deve ser interpretado primordialmente como decorrente de \emph{raciocínio
contextual e correlação entre hospedeiros sob escassez de sinal}, e não como
ganho de revocação bruta sobre eventos inexistentes.

\subsection{Outras ameaças à validade}

\paragraph{Validade interna.} A medição de resiliência opera no nível da
\emph{tentativa} (\textit{trial-level}): a ASR e a NASR medem se o agente
classificou incorretamente sob injeção, mas o corpus não dispõe de rótulos de
contaminação \emph{por evento}. Não se afirma, portanto, qual evento individual
foi corrompido, apenas o desfecho da janela sob contaminação. Adicionalmente, a
regra de \textit{ground truth} baseada no caminho do arquivo
(Equação~\ref{eq:ground-truth}) é determinística e auditável, mas pressupõe que a
curadoria do corpus rotulou corretamente cada \textit{playbook}.

\paragraph{Validade externa.} O ambiente é um domínio \textit{Active Directory}
de pequena escala (um controlador de domínio e duas estações), o que reproduz a
superfície de movimentação lateral de uma pequena empresa, mas não a escala, a
heterogeneidade nem o volume de \textit{background} de uma rede empresarial de
grande porte. O corpus de 185 \textit{playbooks} cobre nove famílias de LOLBin,
porém não esgota o universo de técnicas de \textit{living-off-the-land}, e a
distribuição malicioso/benigno do corpus não pretende refletir a taxa-base
operacional real.

\paragraph{Validade de construto.} A linha de base C2-naive é definida como um
classificador de passagem única; deve-se assegurar, na execução, que ela seja de
fato um passe genuíno de modelo de linguagem (paridade com C1-full), e não uma
heurística determinística, sob pena de a comparação medir algo distinto do
pretendido. Esse cuidado é registrado aqui como ameaça de construto a ser
controlada na fase de execução e auditada nos artefatos (Capítulo~\ref{cap:resultados}).

\paragraph{Estabilidade do laboratório.} A instância QEMU de uma das estações
(ws02) apresentou instabilidade intermitente (asserção do monitor QEMU seguida de
saída do contêiner) logo após a inicialização, sob a pressão de memória do
hospedeiro. A inicialização escalonada, o \textit{watchdog} de reinício e o
re-tentar do gatilho mitigam essa instabilidade na fase de \textit{boot}; uma
falha \emph{durante} a execução, contudo, ainda poderia perder os
\textit{playbooks} roteados para aquele hospedeiro naquela execução. Como
salvaguarda, uma verificação de validade de telemetria, executada após cada
condição, sinaliza qualquer execução com um hospedeiro ausente para
re-execução dirigida.

\bigskip
\noindent
Em síntese, o desenho experimental descrito neste capítulo isola a contribuição
da arquitetura de governança por meio de condições pontuadas sobre janelas e
verdade de referência idênticas, sob um \textit{pipeline} de medição
determinístico e reprodutível, com as suas limitações --- sobretudo a lacuna de
telemetria de EID1 --- documentadas de forma transparente. Os resultados
quantitativos decorrentes deste desenho são apresentados no
Capítulo~\ref{cap:resultados}.
